\chapter{Literature Review}

\section{TOPSIS and Fuzzy TOPSIS}
TOPSIS ranks alternatives by comparing each alternative with a positive ideal solution and a negative ideal solution \cite{hwang1981}. An alternative is preferred when it is closer to the positive ideal and farther from the negative ideal. This structure makes TOPSIS intuitive, computationally tractable, and easy to explain.

In fuzzy TOPSIS, uncertain ratings are often represented as triangular fuzzy numbers. A triangular fuzzy number is denoted as
\[
\tilde{x} = (l,m,u),
\]
where \(l\), \(m\), and \(u\) represent the lower, modal, and upper values. Chen \cite{chen2000} extended TOPSIS for fuzzy group decision-making, making the method suitable for linguistic and uncertain evaluations.

\section{Group Decision-Making}
Group decision-making requires individual judgments to be integrated into a collective decision. A common approach is to aggregate decision-maker ratings using arithmetic or geometric means. However, aggregation can hide disagreement, amplify extreme inputs, or fail when some decision makers are biased.

Huang and Li \cite{huang2012} argued that conventional aggregation of individual TOPSIS judgments can be too intuitive because it may not consider preference priorities and preference intensity. Their group-ideal TOPSIS aggregation model motivates one of the external comparator baselines used in this thesis.

\section{Consensus and Reliability in Group Decisions}
Consensus models aim to measure and improve agreement among decision makers. Liu et al. \cite{liu2022} proposed a consensus-building model for group decision-making with non-reciprocal fuzzy preference relations. Tran et al. \cite{tran2024} reviewed consensus models for group decision-making and group recommender systems, showing that agreement, disagreement, and consensus-achieving processes are central concerns in group decisions.

This thesis uses the consensus idea differently. Instead of forcing decision makers to revise opinions, it estimates reliability from observed rating behavior. Decision makers whose ratings are structurally abnormal or clone-like receive lower influence.

\section{Expert Weighting and Fuzzy MAGDM}
Recent fuzzy multi-attribute group decision-making work also studies expert weighting and optimization-based methods. Hu et al. \cite{hu2023} proposed a fuzzy MAGDM method based on TOPSIS and optimization models to determine expert and attribute weights. This supports the broader idea that decision-maker weights are important, but the present thesis focuses specifically on adversarial or biased decision-maker contamination in triangular fuzzy TOPSIS.

\section{Research Gap}
Prior work has studied fuzzy TOPSIS, group TOPSIS aggregation, consensus models, and expert weighting. However, less attention has been given to targeted rank manipulation caused by coordinated biased decision makers. The key gap is therefore:

\begin{quote}
How can fuzzy TOPSIS group decision-making remain stable when a subset of decision makers intentionally pushes a weak target alternative upward?
\end{quote}

This thesis addresses that gap through reliability-aware ensemble methods and structured contamination experiments.
